% Competitive reversal: the Pimentel model --- plan \S3.3, "\S3".
% Sources: 09/05:1-205.  Heading spelling corrected throughout: the source has
% "Pimental" in every heading and "Pimentel" in the body and the bib key.
% Rename: evolution rate \mu -> \varepsilon (plan \S7.2).
% Repairs: the alpha = 1 to alpha = 10 jump at 09/05:78-83, and the forward
% reference to nodisp2runs at 09/05:80.

\section{Competitive reversal: the Pimentel model}
\label{var:sec:pimentel}

The speciation model of the previous section lets variation out when there
is room. This one asks the complement: what happens when two populations
share the room, and the minority is the one under selection.

\subsection{The experiment}
\label{var:sec:pimentelexperiment}

Pimentel and colleagues devised and ran an experiment in which two
populations of flies are pitted against one another in a fight for domination
and, ultimately, survival \cite{pimentel1965selection,hutchingson2018ecological}.
Place a colony of houseflies (\textit{Musca domestica}) in one side of a
boxed enclosure and a competing number of blowflies (\textit{Phaenicia}) at
the other, and give them an ample but finite supply of food. Before long one
of the species dominates the niche. The arrangement is unstable, and either
the housefly or the blowfly population eventually sees off the other.

Pimentel and his colleagues did something clever to augment the dynamics.
They contrived an enclosure --- the original box, but with a myriad of
internal sub-compartments, a three-dimensional maze of tunnels and small
chambers. That arrangement gives the enclosure a desirable property: with the
populations initially placed at opposing sides, communication by contact is
delayed by a period greater than the average lifespan of either species. In
the months taken for a species to spread from one side to the other, multiple
generations ensue.

Why is that desirable? By slowing the potential spread, Pimentel imposes the
constraint that a total domination, if one is arrived at, cannot be won in
one or a few fly lifespans. The war must be won over succeeding generations,
grandchildren's grandchildren all fighting for the same cause, and that is
what allows the hand of Darwinian selection to come into play.

The observed dynamic was the following. Blowflies were pushed almost to
extinction by the initially superior numbers of \textit{Musca domestica}, and
then the tide turned: from the brink, the blowfly underdogs clawed back their
footing and eventually overthrew the advances of their adversaries. The
hypothesis is that during this squeeze of reduced population the blowflies
began to select, among their limited numbers, for the abilities favouring
direct competition with their adversaries. The houseflies, doing well as they
were, were subject to no such channelling selection; they became complacent
while their foes grew strong, cut into shape by a harshly oppressive
environment.

It has not been achieved experimentally, but Pimentel proposed that with the
enclosure arranged correctly an oscillation could be set up, in which the
upper hand passes continually from one side to the other, each being pushed
to the brink and moulded to competitive advantage in times of scarcity --- a
continual arms race for survival. The question this section asks is whether a
Markov chain with evolving blocking probabilities can reverse in the same way.

\subsection{The model}
\label{var:sec:pimentelmodel}

Consider two populations competing for space in a finite domain they both
occupy. Let $A_t$ count one of them, labelled $a$, and $B_t$ the other,
labelled $b$, both in discrete time $t\in\mathbb{N}\cup\{0\}$. Finiteness of
the enclosure is imposed by $A_t+B_t\le N$ for a maximum capacity $N$, and
with the initial condition $A_0=B_0=N/2$ for even $N$ the two counts are tied
by $B_t=N-A_t$ throughout, so it is enough to follow $A_t$. Its transition
probabilities are
\begin{equation}
\pr{A_{t+1}=j\mid A_t=i}=
\begin{cases}
q_a=\dfrac{i}{N}(1-p_b), & j=i+1,\\[4pt]
q_b=\lrb{1-\dfrac{i}{N}}(1-p_a), & j=i-1,\\[4pt]
q_c=\dfrac{i}{N}p_b+\lrb{1-\dfrac{i}{N}}p_a, & j=i,\\[4pt]
0 & \text{otherwise,}
\end{cases}
\label{var:eq:kernel}
\end{equation}
for $0<i<N$, with $i=0$ and $i=N$ absorbing.

The three cases come from the following step rule, viewed from the
perspective of population $a$.

\begin{description}[leftmargin=1.6em,style=unboxed,itemsep=0.15em]
\item[Whose move.] With probability $i/N$ population $a$ attempts to advance;
with probability $1-i/N$ the opponent attempts to, potentially forcing $a$
back one.
\item[Advancing.] If $a$ has decided to advance it does so, unless it is
blocked, which happens with probability $p_b$.
\item[Being advanced upon.] If the opponent is advancing, $a$ blocks with
probability $p_a$; if it fails to, it is forced back a step.
\item[Stalemate.] If a block is effective, both parties remain where they
are, and $A_{t+1}=A_t$.
\end{description}

The capacity to block is the model's evolving quantity. Let it depend on a
\emph{resistance} $R_k$ for $k=a,b$ through
\begin{equation}
p_a=\max\left\{0,\,1-\frac{R_b}{R_a}\right\},
\qquad
p_b=\max\left\{0,\,1-\frac{R_a}{R_b}\right\},
\label{var:eq:blockprobs}
\end{equation}
so that at most one of the two is non-zero: whichever population has the
higher resistance blocks, and the other does not. If $R_a\gg R_b$ then
$p_a$ is close to one and $p_b$ is zero.

Resistances are driven up when the associated population is in a minority.
They are updated at each step by
\begin{equation}
R_{k,t+1}=R_{k,t}+\varepsilon\,\Delta R_{k,t},
\label{var:eq:resistupdate}
\end{equation}
at an evolution rate $\varepsilon$, with
\begin{equation}
\Delta R_{a,t}=\max\left\{0,\lrb{\frac{N-A_t}{A_t}}^{\!\alpha}\right\},
\qquad
\Delta R_{b,t}=\max\left\{0,\lrb{\frac{A_t}{N-A_t}}^{\!\alpha}\right\}.
\label{var:eq:dR}
\end{equation}
Each increment grows as its population is squeezed towards its absorbing
boundary, and is larger for the population in the minority. Only the ratio
$R_b/R_a$ enters \cref{var:eq:blockprobs}, so what drives the blocking
probabilities is the \emph{difference} between the two increments rather than
either alone. The
mechanism intended is the one Pimentel proposed: most of the competitive
interactions experienced by members of the smaller population are with the
opposition, and the drive for selection towards competitive ability is
correspondingly stronger, while the majority population meets mostly its own
kind and is under no such pressure. The exponent $\alpha$ sharpens that
drive near the boundaries. Taking $\alpha=1$ makes the increment simply
proportional to the ratio of interspecific to intraspecific encounters, which
is the physical case; $\alpha>1$ is retained because it is what makes the
criticality of \cref{var:sec:critical} visible.

The initial resistances need saying, because the model is not scale-free in
them. Since \cref{var:eq:blockprobs} depends on $R_a$ and $R_b$ only through
their ratio, while \cref{var:eq:resistupdate} adds an absolute increment,
scaling $R_{a,0}$, $R_{b,0}$ and $\varepsilon$ by a common factor leaves every
trajectory unchanged: the two enter only through $\varepsilon/R_0$. Taking
$R_{a,0}=R_{b,0}=1$ throughout therefore costs nothing and fixes what
$\varepsilon$ means.

Two things about \cref{var:eq:kernel,var:eq:blockprobs,var:eq:dR} are worth
naming, since they are what make this a chapter about moving rates. The first
is that the step probabilities are functions of the resistances and the
resistances are functions of the state, so the transition kernel at time $t$
depends on the whole path up to $t$ and not on $A_t$ alone. The second follows:
the triple $(A_t,R_{a,t},R_{b,t})$ is Markov, and $A_t$ by itself is not.

\subsection{Results}
\label{var:sec:pimentelresults}

Simulated at $\alpha=1$, the model does what Pimentel observed. Both
realisations in \cref{var:fig:pim-two-runs} take the same shape: one
population reaches a position of real advantage, and does not keep it.
Trained by interspecific competition, the oppressed population makes a
striking comeback and overthrows its adversary completely. At this exponent
the reversal happens once or twice and then the chain is absorbed, so what
Pimentel hoped for arrives as an episode rather than as a cycle.

\begin{figure}[!htbp]
\centering
\includegraphics[width=0.72\textwidth]{fig03_pimentel_runs/fig03}
\caption{Two realisations of the Pimentel chain at $\alpha=1$, $N=100$,
$\varepsilon=10$, the physical case. Each reverses: an early advantage is
not kept, and the two histories, started from the same symmetric state,
absorb at opposite boundaries.}
\label{var:fig:pim-two-runs}
\end{figure}

Raising $\alpha$ sharpens the resistance drive near the boundaries, and with
it the sensitivity of a trajectory to where it starts.
\Cref{var:fig:pim-runs} makes that comparison at $\alpha=10$ and
$N=10^{5}$, running two independent realisations from each of two
initial conditions: one with the populations evenly matched, one with a bias
imposed at the outset. Started well matched, the two realisations diverge
substantially; started biased, they stay close for far longer. What differs
is not the noise, which is the same in both, but where the trajectory sits
relative to the state at which the chain has no preferred direction.

Raising $\alpha$ also buys far more of the behaviour the model was built for.
Where $\alpha=1$ gives one reversal or two, $\alpha=10$ buys of order ten
reversals before fixation, and the swing grows rather than decays.
\Cref{var:fig:pim-runs} under-shows that run: the full excursion, at the same
parameters, continues for several million further steps before absorption.
The oscillation Pimentel conjectured is therefore not absent from the model,
only impermanent. Every run here still fixates in the end, and what raising
$\alpha$ moves is how long the alternation lasts before it does.

\begin{figure}[!htbp]
\centering
\includegraphics[width=0.96\textwidth]{fig04_pimentel_ic/fig04}
\caption{Two independent realisations from each of two initial conditions, at
$\alpha=10$, $N=10^{5}$, $\varepsilon=1$, over $2\times10^{6}$ steps.
Panel~(a): started evenly matched, and so near the critical state, the runs
separate early and stay apart. Panel~(b): started with one in the majority,
and so away from it, they track one another. Proximity to the critical
state, not the size of the noise, decides how far apart two histories of the
same model end up.}
\label{var:fig:pim-runs}
\end{figure}

\subsection{The moving critical point}
\label{var:sec:critical}

That state can be located exactly. The chain is at its \emph{critical point}
when a step up and a step down are equally likely, that is when $q_a=q_b$ in
\cref{var:eq:kernel}. Writing $i_c$ for the state at which this holds and
using the mutual exclusivity of $p_a$ and $p_b$,
\begin{equation}
i_c=
\begin{cases}
N\dfrac{1-p_a}{2-p_a}, & p_a>0 \text{ and } p_b=0,\\[8pt]
N\dfrac{1}{2-p_b}, & p_b>0 \text{ and } p_a=0.
\end{cases}
\label{var:eq:icrit}
\end{equation}
When neither population blocks, $p_a=p_b=0$ and $i_c=N/2$, the symmetric
point. Blocking moves it off centre, and moves it towards the boundary the
\emph{leading} population is heading for: if $a$ is in the majority then it
is $b$ that accumulates resistance, so $p_b>0$ and $i_c=N/(2-p_b)$ climbs
towards $N$. And $i_c$ is an \emph{unstable} point --- comparing $q_a$
with $q_b$ in \cref{var:eq:kernel} gives $q_a>q_b$ exactly when $i>i_c$, in
both cases of \cref{var:eq:icrit} --- so the chain is carried away from it
rather than returned to it. A reversal is what happens when $i_c$, climbing
after the leader, overtakes the trajectory, and the drift that was carrying
the majority forward begins to carry it back.

\Cref{var:eq:icrit} is less analytic than it looks. Since $p_a$ and $p_b$
depend on $R_a$ and $R_b$, and those are accumulated along a stochastic
trajectory by \cref{var:eq:resistupdate}, $i_c$ is a function of the path and
not of the parameters alone. What it is not is intractable. Given a
realisation, the sequence $i_{c,t}$ is computed step by step at the cost of
the simulation itself and can then be drawn over the trajectory that produced
it, which is what \cref{var:fig:pim-critical} does.

\begin{figure}[!htbp]
\centering
\includegraphics[width=0.72\textwidth]{fig05_pimentel_icrit/fig05}
\caption{A single realisation at $\alpha=3$, $N=10^{4}$, $\varepsilon=1$,
with the critical state $i_{c,t}$ of \cref{var:eq:icrit} overlaid. The
critical state does not stay put. It is dragged towards whichever boundary
the leading population is heading for, and saturates near $0$ or near $N$
while that population is well ahead; each reversal of $A_t$ follows $i_{c,t}$
crossing it.}
\label{var:fig:pim-critical}
\end{figure}

Read against \cref{var:fig:pim-runs}, \cref{var:fig:pim-critical} makes the
mechanism plain: trajectories diverge exactly where $A_t$ is near $i_{c,t}$,
because there the drift that would otherwise restore them is absent. That is
what the resistance dynamics buy. An ordinary random walk on $\{0,\dots,N\}$
is critical only if its step probabilities happen to be balanced, whereas
here the resistance update \emph{drags the critical point after the leading
population} and so keeps returning the process to a state at which its own
future is undetermined. The reversal Pimentel observed is therefore not a
deterministic consequence of the squeeze; it is what a squeeze does to a
chain that the squeeze itself has made critical.

\subsection{A spatial version}
\label{var:sec:automaton}

The model above has no geometry, where Pimentel's enclosure had a great deal
of it. A spatial version was therefore written as a cellular automaton, in
\texttt{C++} using the SFML graphics library, each cell holding one of the two
populations and resistances being inherited and updated locally by the same
rule; \cref{var:fig:automaton} shows six frames from it. A model of
evolutionary reversal in competitive dominance had been proposed before
\cite{pease1984evolutionary}; the model here was not derived from it and
differs in form.

\begin{figure}[!htbp]
\centering
\includegraphics[width=0.96\textwidth]{fig06_automaton/fig06}
\caption{Six frames from the spatial version, left to right along the top row
then the bottom: advance, squeeze, fringe, crossing, recovery, fixation.
Unoccupied cells are green in the first frame. Red advances until blue is
confined to one corner, where a pale fringe of elevated resistance appears;
blue then crosses back and takes the domain. The reversal is spatial as well
as numerical.}
\label{var:fig:automaton}
\end{figure}

% Coupling speciation to competition --- plan \S3.3, "\S4".
% Sources: 09/05:206-323, minus the block quotation at :213-290 (parked as
% _parked/nietzsche_bge_262.tex) and minus the confession at :321-323 (cut;
% replaced by the "What this chapter is" paragraph of \S1).
% Renames beyond plan \S7.2, both forced by \S7.1's rule against contradicting
% a thesis-wide symbol: the fitness increment \delta -> \varepsilon (\delta is
% the catastrophe rate, and this quantity is the evolution rate of the
% preceding section under a second letter), and the speciation exponent
% \lambda -> \zeta (\lambda is the birth rate).  See PHASE_A_REPORT.md 9.1.
% The source at 09/05:276,278 prints both step probabilities with the same
% left-hand side, \Pr(\Delta A_t = 1); the second is written here as -1.

\subsection{Coupling speciation to competition}
\label{var:sec:pimentelplus}

\Cref{var:sec:pimentel} makes a losing population strong, the squeeze being
what drives its resistance up. This subsection asks the complement, whether a
winning population can be made weak by a mechanism of the same kind; the
hypothesis it states is that excess variation weakens an established type, so
that a population which is winning diversifies and thereby loses. The claim
is Nietzsche's, put into a model: a type is made hard by long struggle
against constant unfavourable conditions, favourable conditions relax the
constraint which made it hard, and the resulting exuberance of variation
is itself the mechanism of decline \cite{nietzsche2013beyond}. What is worth retaining is the shape of the claim: that the decline is
\emph{endogenous}, a consequence of the success preceding it rather than a
response to any external shock.

The model is the kernel of \cref{var:sec:pimentelmodel}, recast so that the
quantity under selection is fitness rather than resistance, with a species
count carried alongside and allowed to act back on that fitness. Written as
step probabilities for $A_t$,
\begin{equation}
\pr{\Delta A_t=1}=\frac{A_t}{N}\,\frac{f_a}{\bar f},
\qquad
\pr{\Delta A_t=-1}=\lrb{1-\frac{A_t}{N}}\frac{f_b}{\bar f},
\label{var:eq:fitkernel}
\end{equation}
with $\bar f$ the mean fitness across the domain and, under the Moran
assumption, $A_t=N-B_t$. Each probability is again the chance of meeting a
member of the opposing group multiplied by the chance of prevailing against
them, as in \cref{var:eq:kernel}; what has changed is that prevailing is now
governed by a fitness ratio rather than by a blocking probability. Fitness is
driven up by minority status exactly as resistance was:
\begin{equation}
f_a(t+\Delta t)=f_a(t)+\varepsilon\lrb{\frac{B_t}{A_t}}^{\!\alpha},
\qquad
f_b(t+\Delta t)=f_b(t)+\varepsilon\lrb{\frac{A_t}{B_t}}^{\!\alpha}.
\label{var:eq:fitupdate}
\end{equation}

Fitness here is the same evolving quantity as the resistance $R_a,R_b$ of
the previous section, under a second name: the update rule and its exponent
are identical, and the increment coefficient is $\varepsilon$ in both. The
two kernels are not literally identical, since \cref{var:eq:kernel} has a
stalemate outcome and \cref{var:eq:fitkernel} has none; both are carried
because the fitness form is the one that admits the coupling below, in which
a scalar can be added to and subtracted from.

Now give each population its own species count, $S^{(a)}_t$ and $S^{(b)}_t$,
evolving as birth--death processes with constant per-species extinction
$\mu_a$ and a speciation rate that rises with both dominance and fitness:
\begin{equation}
\pr{\Delta S^{(a)}_t=1}=\beta_a S^{(a)}_t\,\Delta t,
\qquad
\pr{\Delta S^{(a)}_t=-1}=\mu_a S^{(a)}_t\,\Delta t,
\qquad
\beta_a=\lrb{\frac{A_t}{N}}^{\!\zeta}f_a,
\label{var:eq:subspec}
\end{equation}
and correspondingly for $b$ with $A_t$ replaced by $B_t$, the exponent
$\zeta$ controlling how firmly dominance must be established before variation
is released. This is the premise of \cref{var:sec:speciation} transplanted
into a competitive setting: variation is in constant supply, and what
releases it is the easing of whatever was suppressing it, measured here by
the share of the domain a population controls.

The coupling that closes the loop is a penalty on fitness proportional to the
species count:
\begin{equation}
f_a(t+\Delta t)=f_a(t)+\varepsilon\lrb{\frac{B_t}{A_t}}^{\!\alpha}
-\gamma S^{(a)}_t,
\qquad
f_b(t+\Delta t)=f_b(t)+\varepsilon\lrb{\frac{A_t}{B_t}}^{\!\alpha}
-\gamma S^{(b)}_t,
\label{var:eq:fitpenalty}
\end{equation}
with $\gamma$ the amount by which the inefficiency of carrying excess
variation weighs on competitive performance.

That is the whole specification, and it has never been simulated. Both routes
to reversal are present in it at once: the squeeze that strengthens the
loser, through \cref{var:eq:fitupdate}, and the variation that weakens the
winner, through the penalty of \cref{var:eq:fitpenalty}. Both favour whichever
population is behind, so the model does not say in advance which of them does
the work, and that is what makes it worth running.

\paragraph{Open questions.}\label{var:rem:open-plus}
\begin{enumerate}[leftmargin=1.6em,itemsep=0.2em,topsep=0.2em]
\item \textbf{Whether collapse by variation happens at all.} Does the
$-\gamma S^{(a)}_t$ term in \cref{var:eq:fitpenalty} ever produce a reversal
that the squeeze-driven mechanism of \cref{var:sec:pimentel} would not have
produced unaided, or does it only slow that one down? Separating the two
means holding everything else fixed and comparing runs at $\gamma=0$ against
runs at $\gamma>0$, which asks for simulation and for nothing else.
\item \textbf{The relevant regime.} Collapse by variation should require
$\gamma S^{(a)}_t$ to grow comparable with the accumulated fitness advantage
over the timescale on which $S^{(a)}_t$ itself grows, which ties $\gamma$,
$\zeta$ and $\mu_a$ together. Whether that regime is a thin sliver of
parameter space or a broad band of it follows directly from the first
question.
\end{enumerate}


\paragraph{Open questions.}\label{var:rem:open-pimentel}
\begin{enumerate}[leftmargin=1.6em,itemsep=0.2em,topsep=0.2em]
\item \textbf{Fixation.} The fixation probability and the expected time to
fixation, as functions of $(\alpha,\varepsilon,N)$, were never computed. The
first-step machinery of \cref{m:prop:firststep} applies to the enlarged chain
$(A_t,R_{a,t},R_{b,t})$, which is where the work would start.
\item \textbf{Whether the oscillation persists.} Every realisation run here
eventually fixated, at $\alpha=1$ after an episode and at $\alpha=10$ after
of order ten reversals. So the question is not whether the alternation exists
but whether it can be made to last: does the expected number of reversals
before fixation grow without bound in $\alpha$, $\varepsilon$ or $N$, or does
it saturate? This is the sharpest question in the section, because it is the
one the model was built to answer.
\item \textbf{A diffusion limit.} For large $N$ there should be a diffusion
approximation in which \cref{var:eq:icrit} becomes a moving unstable
equilibrium and the divergence of nearby trajectories is a statement about a
local Lyapunov exponent. That would turn the qualitative observation of
\cref{var:fig:pim-critical} into a rate.
\end{enumerate}

\FloatBarrier
